% !TeX root = ../report.tex

\section{GIỚI THIỆU}
\label{sec:gioi-thieu}

\subsection{Phát biểu vấn đề}
\label{subsec:phat-bieu-van-de}

Sự cố cháy nổ tại các khu dân cư, tòa nhà cao tầng, phòng thí nghiệm và nhà xưởng sản xuất luôn tiềm ẩn thiệt hại đặc biệt nghiêm trọng về tính mạng con người và tài sản. Theo tổng hợp từ các nghiên cứu về hệ thống phát hiện cháy dựa trên Internet vạn vật (IoT) \cite{ref_fire_survey}, khoảng thời gian vàng để phát hiện đám cháy ở giai đoạn khởi phát (cháy âm ỉ sinh khói và tăng nhiệt cục bộ) trước khi bùng phát toàn diện thường chỉ kéo dài từ vài chục giây đến vài phút. Tuy nhiên, phần lớn các hệ thống báo cháy dân dụng phổ biến hiện nay vẫn bộc lộ bốn hạn chế kỹ thuật lớn:

\begin{enumerate}[leftmargin=*, itemsep=3pt]
  \item \textbf{Tỉ lệ báo động giả cao do phụ thuộc đơn cảm biến:} Các đầu báo khói hoặc nhiệt độc lập dựa trên một ngưỡng cố định dễ kích hoạt nhầm trước hơi nước, bụi mịn hoặc khói nấu nướng, làm suy giảm độ tin cậy của người vận hành.
  \item \textbf{Phụ thuộc đường truyền đám mây và mất dữ liệu khi gián đoạn mạng:} Trong nhiều mô hình IoT kết nối trực tiếp tới máy chủ đám mây, khi mất kết nối Internet (tình huống rất dễ xảy ra ngay khi sự cố cháy làm đứt cáp hoặc mất điện bộ định tuyến), dữ liệu giám sát bị thất thoát hoàn toàn và người vận hành tại chỗ mất khả năng quan sát, điều khiển.
  \item \textbf{Thiếu khả năng phản ứng liên động tại biên:} Nhiều hệ thống chỉ dừng lại ở mức phát còi báo động thụ động mà không kết hợp điều khiển cơ cấu chấp hành (như đóng cửa ngăn khói, kích hoạt đèn chỉ thị thoát hiểm) hoặc thiếu cơ chế khóa liên động an toàn giữa lệnh điều khiển từ xa và trạng thái cảm biến thực tế.
  \item \textbf{Hạn chế về phân tích thông minh theo thời gian thực:} Việc gửi toàn bộ dữ liệu thô lên đám mây rồi chờ phản hồi làm tăng độ trễ và không tận dụng được năng lực tính toán ngay trên vi điều khiển (TinyML) cũng như máy tính biên (Local IoT Gateway).
\end{enumerate}

Xuất phát từ bài toán thực tiễn đó, nhóm nghiên cứu và phát triển \textbf{Hệ thống Báo cháy và Điều khiển Thoát hiểm Thông minh (IoT Fire Alarm System)} dựa trên nền tảng vi điều khiển ESP32-S3 (YoloUNO) \cite{ref_esp32s3}, kết hợp kiến trúc \textbf{Dual-AI} (TinyML chạy trực tiếp trên vi điều khiển song song với thuật toán Isolation Forest \cite{ref_iforest} trên Local IoT Gateway), cơ chế lưu đệm ngoại tuyến SQLite Write-Ahead Logging (WAL) \cite{ref_sqlite_wal} và nền tảng quản trị CoreIOT/ThingsBoard \cite{ref_thingsboard}. Đối tượng sử dụng của hệ thống là ban quản lý tòa nhà, kỹ thuật viên trực phòng thí nghiệm và hộ gia đình cần một giải pháp cảnh báo cháy có độ tin cậy cao, duy trì hoạt động độc lập tại chỗ khi mất mạng (Local Survivability) và hỗ trợ giám sát -- điều khiển từ xa qua Internet.

\subsection{Các giải pháp liên quan}
\label{subsec:giai-phap-lien-quan}

Để xác định khoảng trống công nghệ mà đề tài hướng tới, nhóm tiến hành khảo sát ba hướng tiếp cận phổ biến đối với hệ thống báo cháy hiện nay:

\begin{itemize}[leftmargin=*, itemsep=3pt]
  \item \textbf{Hệ thống báo cháy truyền thống cục bộ (Standalone Fire Alarm):} Gồm các đầu báo khói quang/ion và tủ trung tâm có dây. Ưu điểm là độ ổn định phần cứng cao, không phụ thuộc Internet; nhược điểm là không hỗ trợ giám sát từ xa, không lưu lịch sử chuỗi thời gian phục vụ truy vết và thiếu khả năng phân tích đa biến bằng học máy.
  \item \textbf{Thiết bị báo cháy IoT kết nối trực tiếp Cloud (Direct-to-Cloud IoT):} Sử dụng các vi điều khiển tích hợp Wi-Fi gửi trực tiếp dữ liệu qua giao thức MQTT \cite{ref_mqtt} hoặc HTTP lên nền tảng đám mây \cite{ref_fire_survey}. Ưu điểm là triển khai nhanh, chi phí thấp; tuy nhiên khi đường truyền Internet bị gián đoạn, thiết bị không có bộ đệm lịch sử dung lượng lớn, không có giao diện giám sát cục bộ tại trạm và việc xử lý quy tắc thông minh trên đám mây bị ngắt quãng.
  \item \textbf{Hệ thống báo cháy ứng dụng học máy đám mây hoặc thị giác máy tính:} Triển khai mô hình mạng nơ-ron sâu trên máy chủ trung tâm để nhận diện khói/lửa. Ưu điểm là độ chính xác nhận dạng cao trong môi trường phức tạp; nhược điểm là yêu cầu băng thông lớn, chi phí phần cứng GPU cao và độ trễ phụ thuộc chất lượng mạng.
\end{itemize}

Bảng \ref{tab:so-sanh-giai-phap} tổng hợp so sánh các hướng tiếp cận trên với hệ thống nhóm đề xuất.

\begin{table}[htbp]
  \centering
  \caption{So sánh các giải pháp báo cháy liên quan và hệ thống đề xuất}
  \label{tab:so-sanh-giai-phap}
  {\small
  \setlength{\tabcolsep}{4pt}
  \renewcommand{\arraystretch}{1.22}
  \begin{tabularx}{\textwidth}{|L{0.19\textwidth}|L{0.17\textwidth}|L{0.17\textwidth}|L{0.17\textwidth}|Y|}
    \hline
    \rowcolor{iotTableHead}
    \textbf{Tiêu chí} & \textbf{Báo cháy truyền thống} & \textbf{IoT nối thẳng Cloud} & \textbf{AI trên Cloud/GPU} & \textbf{Hệ thống nhóm đề xuất} \\
    \hline
    Đa cảm biến & Hạn chế & Có & Chủ yếu camera & MQ-2 + Flame + DHT11 \\
    \hline
    AI tại biên (Edge/TinyML) & Không & Không & Chạy trên server & Dual-AI: TinyML (MCU) + Isolation Forest (Gateway) \\
    \hline
    Chịu lỗi mất mạng & Tốt (chỉ tại chỗ) & Kém (mất gói) & Kém & Đệm SQLite WAL + SCADA cục bộ cổng 5050 \\
    \hline
    Điều khiển hai chiều & Thủ công tại tủ & Qua Cloud & Qua Cloud & Tự động RuleChain + RPC từ xa + Local Override \\
    \hline
    Chi phí triển khai & Trung bình & Thấp & Cao & Thấp (ESP32-S3 + PC/Docker local) \\
    \hline
  \end{tabularx}
  }
\end{table}

Khoảng trống mà đề tài giải quyết là kết hợp tính tức thời và độc lập tại chỗ của hệ thống nhúng (suy luận TinyML trực tiếp trên MCU, trễ cục bộ tức thì), khả năng chịu lỗi mất mạng kết hợp phát hiện bất thường đa biến trên Local Gateway (SQLite WAL + Isolation Forest + SCADA Local) và khả năng tự động hóa, quản trị tập trung trên nền tảng CoreIOT/ThingsBoard (Dashboard công nghiệp + Rule Chain + Notification).

\subsection{Mục tiêu đề tài}
\label{subsec:muc-tieu-de-tai}

\textbf{Mục tiêu tổng quát:} Thiết kế, hiện thực và kiểm chứng hoàn chỉnh hệ thống báo cháy IoT đa tầng gồm End-Node ESP32-S3 (YoloUNO), Local IoT Gateway trên PC và máy chủ CoreIOT/ThingsBoard tự triển khai cục bộ, đảm bảo phát hiện cháy chính xác bằng kiến trúc Dual-AI, điều khiển cơ cấu chấp hành tự động/thủ công hai chiều và không thất thoát dữ liệu khi gián đoạn mạng.

\textbf{Các mục tiêu cụ thể, đo lường được:}
\begin{enumerate}[leftmargin=*, itemsep=3pt]
  \item \textbf{Thu thập đa cảm biến và phản ứng tại chỗ:} Node ESP32-S3 lấy mẫu cảm biến khói MQ-2 (ADC 12-bit) và cảm biến lửa với chu kỳ $100\text{ ms}$, đọc nhiệt độ -- độ ẩm DHT11 mỗi $2000\text{ ms}$, cập nhật màn hình TFT ILI9341 $320 \times 240$ và truyền telemetry qua UART $115200\text{ bps}$ mỗi $500\text{ ms}$.
  \item \textbf{Suy luận Dual-AI chính xác và nhẹ:} Xây dựng công cụ huấn luyện và xuất tự động mô hình TinyML sang tệp tiêu đề C (\texttt{tinyml\_fire\_model.h}) đạt độ chính xác trên tập kiểm thử $> 95\%$, tiêu thụ RAM $< 1\text{ KB}$ trên ESP32-S3; đồng thời triển khai bộ phát hiện bất thường Isolation Forest trên Local Gateway phân tích vectơ 3 chiều (khói, nhiệt độ, lửa).
  \item \textbf{Độ trễ điều khiển vòng lặp thấp:} Thời gian phản hồi vòng lặp đầu--cuối (End-to-End Round-Trip) từ thời điểm gửi bản tin telemetry cảnh báo cháy qua MQTT QoS 1 tới máy chủ ThingsBoard đến khi nhận đủ cả 3 lệnh điều khiển RPC (\texttt{setServo}, \texttt{setBuzzer}, \texttt{setLed}) từ Rule Chain đạt $< 500\text{ ms}$.
  \item \textbf{Độ tin cậy đồng bộ và lưu đệm ngoại tuyến:} Khi kết nối giữa Gateway và Broker bị gián đoạn, $100\%$ bản tin telemetry được lưu đệm an toàn vào cơ sở dữ liệu SQLite WAL và tự động đồng bộ bù theo thứ tự FIFO ngay khi kết nối phục hồi (tỉ lệ mất gói cục bộ $0\%$).
  \item \textbf{Tự triển khai hạ tầng và kiểm thử tự động:} Tự triển khai máy chủ ThingsBoard CE trên Docker cục bộ (cổng HTTP 9090, MQTT 1883), cung cấp đầy đủ bộ script tự động cấp phát Dashboard 10 widget, Rule Chain và bộ kiểm thử đơn vị tự động.
\end{enumerate}

\subsection{Đề xuất tính năng hệ thống}
\label{subsec:de-xuat-tinh-nang}

Toàn bộ yêu cầu chức năng của hệ thống được định danh từ \texttt{F01} đến \texttt{F10} trong Bảng \ref{tab:danh-sach-tinh-nang} và được sử dụng để ánh xạ xuyên suốt các chương sau.

{\small
\setlength{\tabcolsep}{4pt}
\renewcommand{\arraystretch}{1.22}
\begin{longtable}{|C{0.07\textwidth}|L{0.51\textwidth}|L{0.22\textwidth}|C{0.11\textwidth}|}
  \caption{Danh sách tính năng của hệ thống}
  \label{tab:danh-sach-tinh-nang} \\
  \hline
  \rowcolor{iotTableHead}
  \textbf{Mã} & \textbf{Mô tả tính năng} & \textbf{Thiết bị thực hiện} & \textbf{Mức ưu tiên} \\
  \hline
  \endfirsthead
  \hline
  \rowcolor{iotTableHead}
  \textbf{Mã} & \textbf{Mô tả tính năng} & \textbf{Thiết bị thực hiện} & \textbf{Mức ưu tiên} \\
  \hline
  \endhead
  \hline
  \endfoot
  F01 & Thu thập đa cảm biến: MQ-2 (ADC 12-bit, ngưỡng cảnh báo $\ge 1800$, ngưỡng xóa $\le 1650$), cảm biến lửa (active-LOW) mỗi $100\text{ ms}$; cảm biến DHT11 (ngưỡng cảnh báo $\ge 35.0^\circ\text{C}$, ngưỡng xóa $\le 34.0^\circ\text{C}$) mỗi $2000\text{ ms}$; hiển thị trực tiếp lên màn hình TFT ILI9341 $320 \times 240$. & End-Node YoloUNO (ESP32-S3) & Bắt buộc \\
  \hline
  F02 & Suy luận xác suất rủi ro cháy thời gian thực trực tiếp trên vi điều khiển bằng mô hình TinyML (\path{tinyml_fire_model.h}) từ 4 đặc trưng đầu vào và phân loại 3 mức \texttt{SAFE} ($< 40\%$), \texttt{WARNING} ($40\%$--$70\%$), \texttt{CRITICAL} ($\ge 70\%$). & End-Node YoloUNO (ESP32-S3) & Bắt buộc \\
  \hline
  F03 & Công cụ Python tổng hợp dữ liệu vật lý 2.500 mẫu, chuẩn hóa Z-score, huấn luyện bộ phân loại Logistic Regression, đánh giá ma trận nhầm lẫn và tự động xuất mã nguồn C header độc lập. & PC (\path{train_tinyml.py}) & Bắt buộc \\
  \hline
  F04 & Điều khiển cơ cấu chấp hành tại chỗ: đóng/mở cửa ngăn khói bằng động cơ Servo (góc vật lý $0^\circ$/$180^\circ$, quy đổi tỷ lệ $2\times$ từ góc logic $0^\circ$--$90^\circ$), phát còi Buzzer $500\text{ Hz}$ và bật đèn LED cảnh báo; hỗ trợ ghi đè thủ công và ưu tiên ngắt an toàn khi phát hiện cháy thật. & End-Node YoloUNO (ESP32-S3) & Bắt buộc \\
  \hline
  F05 & Cầu nối Local IoT Gateway nhận telemetry từ ESP32-S3 qua UART $115200\text{ bps}$, lưu đệm SQLite WAL chống mất gói, đồng bộ lên CoreIOT/ThingsBoard qua MQTT QoS 1 (hỗ trợ cả Device API và Gateway API) và chuyển tiếp lệnh RPC hai chiều xuống mạch. & Local Gateway (PC) & Bắt buộc \\
  \hline
  F06 & Phát hiện bất thường đa biến thời gian thực tại biên bằng thuật toán học máy Isolation Forest (\texttt{n\_estimators=100}) trên vectơ \texttt{[smoke, temperature, flame]}, tính chỉ số \path{anomaly_score} ($0.00$--$1.00$), cờ \path{is_anomaly} và tự huấn luyện lại trên cửa sổ trượt 200 mẫu sạch. & Local Gateway (PC) & Bắt buộc \\
  \hline
  F07 & Giao diện SCADA Web Monitor cục bộ tại cổng \texttt{5050} hiển thị 5 thẻ chỉ số theo thời gian thực (tần suất cập nhật $1\text{ s}$), bảng điều khiển cưỡng bức cơ cấu chấp hành độc lập với Internet và bảng lưu vết sự kiện thời gian thực (FIFO 25 bản ghi). & Local Gateway (PC) & Bắt buộc \\
  \hline
  F08 & Tự động triển khai và kiểm tra trạng thái máy chủ ThingsBoard CE (\texttt{tb-postgres}) cục bộ qua Docker Compose trên cổng HTTP \texttt{9090} và MQTT \texttt{1883}. & PC (Docker Localhost) & Mở rộng \\
  \hline
  F09 & Tự động cấp phát thiết bị \path{FireAlarm_Gateway} (cờ \texttt{gateway=true}) và triển khai Dashboard công nghiệp 10 widget trên lưới 24 cột (đồng hồ khói, trạng thái lửa, nhiệt độ, AI Anomaly, biểu đồ lịch sử, 3 công tắc RPC, bảng thông báo) trên ThingsBoard/CoreIOT. & CoreIOT / ThingsBoard & Bắt buộc \\
  \hline
  F10 & Xử lý sự kiện cháy tự động bằng Rule Chain trên CoreIOT/ThingsBoard: lọc telemetry, đánh giá nguy cơ cháy, tạo/xóa Alarm mức \texttt{CRITICAL}, kích hoạt thông báo In-App Notification trên Dashboard và điều phối lệnh RPC xuống Gateway. & CoreIOT / ThingsBoard & Bắt buộc \\
  \hline
\end{longtable}
}

\subsection{Phạm vi và giới hạn}
\label{subsec:pham-vi-gioi-han}

Trong khuôn khổ Bài tập lớn môn Phát triển ứng dụng IoT (CO3054), phạm vi hiện thực và giới hạn kỹ thuật của hệ thống được xác định rõ như sau:
\begin{itemize}[leftmargin=*, itemsep=3pt]
  \item \textbf{Phạm vi thực hiện:} Triển khai đầy đủ 1 node phần cứng thực tế ESP32-S3 (YoloUNO) kết nối đầy đủ 3 cảm biến (MQ-2, Flame, DHT11) và 4 ngoại vi hiển thị/chấp hành (TFT ILI9341, Servo, Buzzer, LED); Local Gateway chạy trên PC kết nối mạch qua UART USB-CDC; triển khai ThingsBoard CE cục bộ bằng Docker (và hỗ trợ chuyển đổi sang đám mây \texttt{app.coreiot.io} qua tham số dòng lệnh \texttt{--broker}).
  \item \textbf{Giới hạn kỹ thuật:} Hệ thống được thử nghiệm trong môi trường phòng thí nghiệm với nguồn cấp USB/adapter ổn định, chưa tích hợp mạch pin dự phòng (UPS) chống mất điện khi cháy; tập dữ liệu huấn luyện TinyML được tổng hợp theo đặc tính vật lý cảm biến thay vì thu thập dài ngày từ nhiều công trình thực tế; các tiêu chí không dây giữa nhiều node (ESP-NOW, BLE), Webserver nhúng trực tiếp trên ESP32 (T07), giao diện LVGL (N03) và cập nhật firmware OTA (T14) chưa được kích hoạt trong bản triển khai chính thức này.
\end{itemize}

\subsection{Cấu trúc báo cáo}
\label{subsec:cau-truc-bao-cao}

Phần còn lại của báo cáo được tổ chức thành năm chương: Chương \ref{sec:phan-tich-he-thong} trình bày kiến trúc tổng thể ba tầng, phân rã chức năng, phân lớp công nghệ và bảng ánh xạ tới thang điểm BTL; Chương \ref{sec:hien-thuc-he-thong} mô tả chi tiết quá trình hiện thực từ sơ đồ đấu nối phần cứng, firmware ESP32-S3, Local IoT Gateway, hai mô hình AI (TinyML và Isolation Forest) đến cấu hình ThingsBoard/CoreIOT cùng các minh chứng hình ảnh thực tế; Chương \ref{sec:do-dac-thu-nghiem} trình bày kịch bản thử nghiệm, số liệu đo đạc thực nghiệm (độ trễ vòng lặp E2E, độ chính xác AI, kiểm thử bộ đệm ngoại tuyến) và đánh giá kỹ thuật; Chương \ref{sec:ket-luan-mo-rong} tổng kết kết quả đạt được, phân tích hạn chế và đề xuất hướng phát triển; cuối cùng, Chương \ref{sec:huong-dan-ma-nguon} cung cấp hướng dẫn chi tiết cài đặt môi trường, cấu hình và vận hành mã nguồn trên GitHub.
